\subsection{Transistor-Verbundverstärker}
% Alt: Klausuraufgabe 2
\Aufgabe{
      \label{Aufg27}
      Zwei npn-Bipolartransistoren bilden zusammen einen rauscharmen Verstärker für Signale im Bereich unter $1\, \mathrm{\mu V}$.
      Die Versorgungsspannung beträgt $U_1=10\,\mathrm{V}$. Am Kollektor des Transistors $\mathrm{T_2}$ liegt eine Spannung von $4\,\mathrm{V}$ gegenüber Masse an.
      Im Arbeitspunkt von $T_1$ und $T_2$ kann für die Basis-Emitter-Spannungen jeweils $U_\mathrm{BE}=600\,\mathrm{mV}$ angenommen werden.
      Die Kollektorströme $\mathrm{I_{C}}$ sind sehr groß gegenüber den Basisströmen $I_\mathrm{B}$ anzunehmen, daher gilt näherungsweise $I_\mathrm{C} = I_\mathrm{E}$.
      Zusätzlich sind die folgenden Parameter gegeben:

      $R_1 = 10\, \mathrm{k\Omega}$,
      $R_2 = 10\, \mathrm{k\Omega}$,
      $R_3 = 2\, \mathrm{k\Omega}$,
      $R_4 = 300\, \mathrm{k\Omega}$,
      $R_5 = 10\, \mathrm{M\Omega}$ \\
      $C_1 = 1\, \mathrm{\mu F}$,
      $C_2 = 1\, \mathrm{\mu F}$ \\
      $U_1 = 10\, \mathrm{V}$,
      $\hat{u}_\mathrm{E} = 0,1\, \mathrm{mV}$ \\

      \begin{figure}[H]
            \centering
            \input{Tikz/src/FigKlausurVerbundverstaerker}\alttext{Das Bild zeigt ein elektrisches Schaltbild mit zwei Bipolartransistoren (\(T_1\) und \(T_2\)), fünf Widerständen (\(R_1\) bis \(R_5\)), zwei Kondensatoren (\(C_1\) und \(C_2\)) sowie einer Wechselspannungsquelle \(u_\mathrm{E}(t)\) links und einer Gleichspannungsquelle \(U_1\) oben rechts. Die Bauteile sind durch Leitungen verbunden: Der Wechselspannungseingang führt über \(C_1\) zum Basisanschluss von \(T_1\), dessen Emitter mit Masse verbunden ist. Der Kollektor von \(T_1\) ist über \(R_1\) an die positive Versorgungsspannung angeschlossen und gleichzeitig mit der Basis von \(T_2\) verbunden. Der Emitter von \(T_2\) ist über \(R_3\) mit Masse verbunden, sein Kollektor über \(R_2\) zur Versorgungsspannung. Zwischen Basis von \(T_1\) und Emitter von \(T_2\) liegt \(R_4\). Vom Kollektor von \(T_2\) führt ein Anschluss über \(C_2\) zu \(R_5\), der gegen Masse geschaltet ist. Unten links ist Masse als Bezugspunkt dargestellt.}
            \label{fig:FigKlausurVerbundverstaerker}
      \end{figure}

      \begin{itemize}
            \item[a)]
                  Wie groß ist der maximale Strom durch $R_2$ im Arbeitspunkt ($I_\mathrm{R2}=I_\mathrm{C,\,T2}$)?
            \item[b)]
                  Wie groß ist der Strom durch $R_3$ näherungsweise, wenn $R_4 >> R_3$?
            \item[c)]
                  Wie groß ist die Spannung, die über $R_3$ abfällt?
            \item[d)]
                  Wie groß ist die Gleichspannung an der Basis von $T_2$ ($U_\mathrm{B,\,T2}=U_\mathrm{CE,\,T1}$)?
            \item[e)]
                  Wie groß ist der Strom durch $R_1$ ($I_\mathrm{R1}=I_\mathrm{C,\,T1}$)?
            \item[f)]
                  Erläutern Sie die Arbeitspunktstabiliserung der gesamten Schaltung auf Grund der Gegenkopplung über $R_4$.
            \item[g)]
                  Wie groß ist die Steilheit $S$ von $T_\mathrm{1}$, wenn zusätzlich $U_\mathrm{T} = 26\, \mathrm{mV}$ gegeben ist?
            \item[h)]
                  Wie groß ist die Spannungsverstärkung $A_\mathrm{V} = \frac{\Delta U_\mathrm{CE}}{\Delta U_\mathrm{BE}}$, der ersten Transistorstufe?
            \item[i)]
                  Wie groß ist die Spannungsverstärkung $A_\mathrm{V}$ der zweiten Transistorstufe, unter Verwendung der Definition $A_\mathrm{V,T2}=\frac{\Delta U_\mathrm{C}}{\Delta U_\mathrm{B}}$?
                  Zur Lösung müssen zusätzlich die Widerstände $R_2$ und $R_3$ berücksichtigt werden.
            \item[j)]
                  Wie groß ist die Spannungsverstärkung der beiden Stufen $T_\mathrm{1}$ und $T_\mathrm{2}$ zusammen?
            \item[k)]
                  Zeichnen Sie das Kleinsignal-Ersatzschaltbild der gesamten Schaltung. Ersetzen Sie darin die beiden Transistoren im Arbeitspunkt jeweils durch eine steuerbare Stromquelle $S=\beta\cdot i_\mathrm{B}$ und einen differentiellen Eingangswiderstand $r_\mathrm{BE}$.
      \end{itemize}
}


\Loesung{
      \begin{itemize}
            \item[a)]
                  \begin{align*}
                        I_\mathrm{C} & = \frac{6\,\mathrm{V}}{10\,\mathrm{k\Omega}} = 0,6\,\mathrm{mA}
                  \end{align*}

            \item[b)]
                  \begin{align*}
                        I_\mathrm{C,T2} & = 0,6\,\mathrm{mA}
                  \end{align*}

            \item[c)]
                  \begin{align*}
                        U_{\mathrm{R3}} & = 0,6\,\mathrm{mA} \cdot 2\,\mathrm{k\Omega} = 1,2\, \mathrm{V}
                  \end{align*}

            \item[d)]
                  \begin{align*}
                        1,2\, \mathrm{V} + 0,6\,\mathrm{V} = 1,8\,\mathrm{V}
                  \end{align*}

            \item[e)]
                  \begin{align*}
                        I_\mathrm{C,T1} & = \frac{8,2\,\mathrm{V}}{10\,\mathrm{k\Omega}} = 0,82\,\mathrm{mA}
                  \end{align*}

            \item[f)]
                  Eine Erhöhung von $I_\mathrm{C}$ führt zu einem Spannungsanstieg über $R_4$, dadurch sinkt die Basisspannung von $T_\mathrm{1}$, was $I_\mathrm{C}$ wieder reduziert.
                  Die negative Rückkopplung stabilisiert den Arbeitspunkt.

            \item[g)]
                  \begin{align*}
                        S & = \frac{I_\mathrm{C,T1}}{U_\mathrm{T}} = \frac{0,82\,\mathrm{mA}}{26\,\mathrm{mA}} = 31,54\, \frac{\mathrm{mA}}{\mathrm{V}}
                  \end{align*}

            \item[h)]
                  \begin{align*}
                        A_\mathrm{V} & = S \cdot (-R_\mathrm{C}) = -315
                  \end{align*}

            \item[i)]
                  \begin{align*}
                        A_\mathrm{V} & = -\frac{R_\mathrm{C}}{R_\mathrm{E}} = \frac{10\,\mathrm{k\Omega}}{2\,\mathrm{k\Omega}} = -5
                  \end{align*}

            \item[j)]
                  \begin{align*}
                        -5 \cdot -315 = 1575 \approx 1600
                  \end{align*}

            \item[k)] Ersatzschaltbild\\
                  \begin{minipage}{\linewidth}
                        \begin{figure}[H]
                              \centering
                              \input{Tikz/src/LsgKlausurVerbundverstaerker}
                              \alttext{Die Abbildung zeigt das Kleinsignalersatzschaltbild einer zweistufigen Transistorschaltung. Dargestellt sind zwei hintereinandergeschaltete Transistorstufen mit ihren jeweiligen Ersatzwiderständen und gesteuerten Stromquellen.
                              Links befindet sich die Eingangsspannung \(U_\mathrm{e}\) die zwischen dem linken oberen Knoten und Masse eingezeichnet ist. Dieser Knoten ist mit der Basis des ersten Transistors verbunden und mit „B“ gekennzeichnet. Zwischen Basis und Emitter des ersten Transistors liegt der Basis-Emitter-Widerstand \(r_{\mathrm{BE1}}\) Der Emitter ist mit Masse verbunden.
                              Zwischen Kollektor und Emitter des ersten Transistors ist der Widerstand \(r_{\mathrm{CE1}}\) eingezeichnet. Zusätzlich ist eine gesteuerte Stromquelle dargestellt, deren Strom proportional zum Basisstrom ist und mit \(\mathrm{B} \cdot i_{\mathrm{B1}}\) gekennzeichnet ist. Am Kollektorknoten des ersten Transistors ist zudem der Widerstand \(R_1\) angeschlossen, der zum Emitterknoten führt.
                              Der Kollektorknoten des ersten Transistors ist direkt mit der Basis des zweiten Transistors verbunden. Auch hier ist zwischen Basis und Emitter der Widerstand \(r_{\mathrm{BE2}}\) eingezeichnet. Der Emitter des zweiten Transistors ist über ein Widerstandsnetzwerk mit Masse verbunden. Dieses Netzwerk besteht aus \(R_2\), \(R_3\) und \(R_5\). Zusätzlich ist ein Widerstand \(R_4\)vom Eingangsknoten zu einem Emitterknoten geführt.
                              Zwischen Kollektor und Emitter des zweiten Transistors befindet sich der Widerstand \(r_{\mathrm{CE2}}\). Auch hier ist eine gesteuerte Stromquelle eingezeichnet, deren Strom mit \(\mathrm{B} \cdot i_{\mathrm{B2}}\) bezeichnet ist.
                              Ganz rechts ist die Ausgangsspannung \(u_\mathrm{a}\) eingezeichnet. Sie wird zwischen dem rechten oberen Knoten und dem rechten unteren Knoten eingezeichnet.}       
                        \end{figure}
                  \end{minipage}
      \end{itemize}

      Siehe: Abschnitt \ref{sec:DarlingtonTransistor}
}